\begin{scanpage}{0144}{126}
agrees, when $\sqrt{ka}\cos(\pi/4\pm\varphi/2)\gg1$, with the asymptotic solution obtained in the book [50] by means of integral equations. The solution found in [50] has the greatest precision when $\alpha\approx0$, $\varphi\approx0$, and it is completely useless if $\alpha\approx\pm\pi/2$ or $\varphi\approx\pm\pi/2$.
\sectiontitle{22}{Secondary Diffraction by a Strip (E-Polarization)}
It is known that a current filament with an electric moment $p_z$ which is found at a distance $h$ from an infinite, ideally conducting plane (see Figure 46) creates in space the field
\[
E_z=ik^2p_z\pi[H_0^{(1)}(kR_1)-H_0^{(1)}(kR_2)].\tag{22.01}
\]
With small values of $h$ (and $R_{1,2}\gg kh^2$), this expression is transformed to the form
\[
E_z=-i2p_zk^3h\sin\psi\sqrt{\frac{2\pi}{kR_0}}e^{i(kR_0+\pi/4)}.\tag{22.02}
\]
The primary edge wave is determined by the relationship
\[
E_z=E_{0z}(q)f(q)\frac{e^{i(kR_0+\pi/4)}}{\sqrt{2\pi kR_0}},\tag{22.03}
\]
where $E_{0z}(q)$ is the value of the incident plane wave field at the point $q$ ($R_0=0$). Consequently, the primary edge wave in the direction $\psi=0$ may be investigated as the wave from a current filament located above an ideally conducting plane if one assumes the filament moment to be equal to
\[
p_z=\frac{i}{4\pi k^3h}E_{0z}(q)\frac{f(q)}{\sin\psi}\bigg|_{\psi=0}.\tag{22.04}
\]
The field, created at the point P by the current filament with a moment $p_z$ which is parallel to the half-plane's edge and passes through the point Q, is determined by Expression (20.11). Expanding
\end{scanpage}
